\documentclass[12pt, a4paper]{article}

\usepackage[utf8]{inputenc}
\usepackage[T1]{fontenc}
\usepackage[brazil]{babel}
\usepackage[top=2.2cm, bottom=2.2cm, left=2.8cm, right=2.5cm]{geometry}
\usepackage{booktabs}
\usepackage{multirow}
\usepackage{array}
\usepackage[table]{xcolor}
\usepackage{caption}
\usepackage{float}
\usepackage{microtype}
\usepackage{setspace}
\setstretch{1.04}
\captionsetup{font=small, labelfont=bf, skip=4pt}
\setlength{\abovecaptionskip}{4pt}
\setlength{\belowcaptionskip}{1pt}
\setlength{\parskip}{3pt}

% Cores
\definecolor{azul}    {RGB}{20,  60, 120}
\definecolor{azulhd}  {RGB}{245,249,255}
\definecolor{azulhd2} {RGB}{230,238,252}
\definecolor{verdebg} {RGB}{220,240,220}
\definecolor{verdedark}{RGB}{34, 120, 34}
\definecolor{vermbg}  {RGB}{255,220,220}
\definecolor{vermdark}{RGB}{160, 30, 30}

% Células de matriz de confusão
\newcommand{\cmTP}[1]{\cellcolor{verdebg}\color{verdedark}\textbf{TP\,=\,#1}}
\newcommand{\cmTN}[1]{\cellcolor{verdebg}\color{verdedark}\textbf{TN\,=\,#1}}
\newcommand{\cmFP}[1]{\cellcolor{vermbg}\color{vermdark}\textbf{FP\,=\,#1}}
\newcommand{\cmFN}[1]{\cellcolor{vermbg}\color{vermdark}\textbf{FN\,=\,#1}}

\newcolumntype{C}[1]{>{\centering\arraybackslash}p{#1}}

\usepackage{url}
\usepackage[colorlinks=true, linkcolor=azul, citecolor=azul, urlcolor=azul]{hyperref}

% Compacta espaçamento da lista de referências
\makeatletter
\renewcommand\@openbib@code{%
  \itemsep 1pt \parsep 0pt \topsep 3pt \partopsep 0pt}
\makeatother

\begin{document}

\title{\textbf{Detecção de Falha de Motor em VANTs com Aprendizado de Máquina}\\[4pt]
       \large Resumo de Resultados}
\author{}
\date{}
\maketitle
\thispagestyle{empty}
\vspace{-10pt}

\begin{abstract}
\noindent
Este trabalho apresenta um sistema de detecção automática de falha de motor em voo para a aeronave não tripulada de asa fixa CarbonZ, com base em aprendizado de máquina não supervisionado. Dois modelos \textit{Isolation Forest} foram treinados com trilhas de \textit{features} distintas --- grandezas físicas (\texttt{all\_features}) e domínio da frequência (\texttt{fft\_features}) --- sobre 33 voos em formato ROS (${\sim}520$\,Hz). O modelo físico detectou 16/23 falhas com latência mediana de 3,15\,s e precisão de 0,983; o modelo espectral detectou 15/23 falhas com latência mediana de 2,64\,s, atuando como \textit{precursor} da falha ao custo de mais alarmes falsos.

\smallskip
\noindent\textbf{Palavras-chave}: Detecção de Anomalia, VANT, Falha de Motor, \textit{Isolation Forest}, FFT, \textit{Isolation Forest}.
\end{abstract}

\section{Introdução}

Em voo autônomo, identificar a perda de tração com latência mínima viabiliza resposta de emergência eficaz (pouso de emergência ou autorrotação). O conjunto de dados conta com 33~voos registrados em formato ROS: 23 com falhas induzidas de motor e 10 voos normais, com dados de múltiplos sensores coletados a ${\sim}520$\,Hz \cite{keipour2019automatic}.

\section{Metodologia}

A solução foi construída como um \textit{pipeline} Kedro completo \cite{kedro2020}, da ingestão dos tópicos ROS (alinhamento temporal via \texttt{merge\_asof}) ao treinamento. Duas trilhas paralelas de engenharia de \textit{features} foram desenvolvidas: (i)~\texttt{all\_features} --- grandezas físicas como energia específica, esforço de controle e razão de planeio, acrescidas de estatísticas em janelas deslizantes; (ii)~\texttt{fft\_features} --- potência de pico, entropia espectral e razão de alta frequência dos sinais IMU e magnetômetro. As 20 melhores \textit{features} foram selecionadas pelo $d$~de~Cohen \cite{cohen1988statistical}. O detector escolhido foi o \textit{Isolation Forest} \cite{liu2008isolation} com janelas deslizantes de 20 amostras, regime não supervisionado, adequado à escassez de exemplos rotulados de falha.

\section{Resultados}

A Tabela~\ref{tab:metricas} apresenta as métricas de classificação calculadas com a função \texttt{classification\_report} sobre o conjunto de teste (últimos 30\% de cada voo, em ordem temporal), no nível de janela de 20 amostras.

\begin{table}[H]
  \centering
  \caption{Métricas por modelo e classe (nível de janela de 20 amostras).}
  \label{tab:metricas}
  \small
  \begin{tabular}{llccc}
    \toprule
    \textbf{Modelo} & \textbf{Classe} & \textbf{Precisão} & \textbf{Recall} & \textbf{F1} \\
    \midrule
    \rowcolor{azulhd}
    \texttt{all\_features} & Normal (0) & 0{,}874 & 0{,}999 & 0{,}932 \\
    \rowcolor{azulhd}
    \texttt{all\_features} & Falha  (1) & 0{,}983 & 0{,}184 & 0{,}310 \\
    \rowcolor{azulhd2}
    \texttt{fft\_features} & Normal (0) & 0{,}856 & 0{,}955 & 0{,}903 \\
    \rowcolor{azulhd2}
    \texttt{fft\_features} & Falha  (1) & 0{,}402 & 0{,}158 & 0{,}227 \\
    \bottomrule
  \end{tabular}
\end{table}

O elevado \textit{recall} da classe Normal no \texttt{all\_features} (0{,}999) reflete o comportamento conservador do \textit{Isolation Forest} configurado com \texttt{contamination\,=\,0{,}03}: apenas 53 das 93.289 janelas normais são classificadas como anomalia, correspondendo a uma taxa de falso positivo por janela de apenas 0{,}057\%. Essa conservatividade se inverte na classe de falha --- o \textit{recall} de 0{,}184 indica que o modelo flagra somente as janelas onde os efeitos da falha já se manifestaram de forma inequívoca nas grandezas físicas. As janelas que ele detecta, porém, são quase todas reais: precisão de 0{,}983.

O \texttt{fft\_features} apresenta comportamento distinto. Seu \textit{recall} para a classe Normal é menor (0{,}955), o que se traduz em volume muito maior de alarmes falsos por janela --- em alguns voos, centenas a mais de mil janelas são marcadas como anomalia \textit{antes} do momento oficial da falha. Esse padrão não é ruído aleatório: ele sugere que o modelo captura a degradação espectral progressiva do motor como precursor físico da falha, e não apenas suas consequências. O custo operacional é a menor confiabilidade por alarme (precisão de 0{,}402 na classe de falha).

A Tabela~\ref{tab:voo} resume os indicadores de desempenho no nível de voo, e a Tabela~\ref{tab:cm} apresenta as matrizes de confusão correspondentes. Um voo é classificado como TP se o modelo emitiu ao menos um alarme durante o intervalo de falha, FN caso contrário; voos normais com ao menos um alarme figuram como FP.

\begin{table}[H]
  \centering
  \caption{Resumo de desempenho no nível de voo.}
  \label{tab:voo}
  \small
  \begin{tabular}{lcc}
    \toprule
    \textbf{Indicador} & \textbf{\texttt{all\_features}} & \textbf{\texttt{fft\_features}} \\
    \midrule
    Falhas detectadas     & 16/23 (69{,}6\,\%) & 15/23 (65{,}2\,\%) \\
    Voos normais com FP   & 6/10 (60{,}0\,\%)  & 3/10 (30{,}0\,\%)  \\
    Latência mediana      & 3{,}15\,s          & 2{,}64\,s          \\
    Latência média        & 6{,}72\,s          & 4{,}57\,s          \\
    \bottomrule
  \end{tabular}
\end{table}

\begin{table}[H]
  \centering
  \caption{Matrizes de confusão --- nível de voo.}
  \label{tab:cm}
  \small
  \begin{tabular}{lcc@{\hspace{2em}}lcc}
    \toprule
    \multicolumn{3}{c}{\textbf{\texttt{all\_features}}}
    & \multicolumn{3}{c}{\textbf{\texttt{fft\_features}}} \\
    \cmidrule(r){1-3}\cmidrule(l){4-6}
    & \textbf{Alarme} & \textbf{Silêncio}
    & & \textbf{Alarme} & \textbf{Silêncio} \\
    \textbf{Com falha} & \cmTP{16} & \cmFN{7}
    & \textbf{Com falha} & \cmTP{15} & \cmFN{8} \\
    \textbf{Sem falha}  & \cmFP{6}  & \cmTN{4}
    & \textbf{Sem falha} & \cmFP{3}  & \cmTN{7} \\
    \bottomrule
  \end{tabular}
\end{table}

\section{Discussão}

Os dois modelos capturam fenômenos físicos distintos com perfis de erro opostos. O \texttt{all\_features} reage às \textit{consequências} da falha (queda de energia, acúmulo de erros de controle), sendo conservador: cada alarme tem 98{,}3\,\% de chance de ser genuíno, mas o modelo silencia quando os efeitos ainda não se consolidaram. O \texttt{fft\_features} detecta a assinatura espectral da degradação \textit{antes} do momento oficial da falha; esse comportamento de precursor é evidenciado pelas centenas a milhares de janelas de alarme antecedendo a falha em alguns voos (e.g., 1.455 janelas no voo 2018-10-18~\#3), mas compromete a precisão por alarme (0{,}402) e inviabiliza seu uso isolado.

Cabe destacar que o conjunto de dados é pequeno (33 voos) e o hiperparâmetro \texttt{contamination\,=\,0{,}03} foi escolhido de forma conservadora; um ajuste fino por validação cruzada temporal poderia melhorar o \textit{recall} da classe de falha sem sacrificar a especificidade. Ainda assim, mesmo na configuração atual, um sistema em cascata de dois estágios poderia ser útil: o \texttt{fft\_features} atuaria como detector precoce de baixo limiar e, ao emitir alarme, iniciaria uma janela de confirmação; o \texttt{all\_features} verificaria se os efeitos físicos se manifestaram dentro dessa janela antes de acionar a resposta de emergência. Essa arquitetura preservaria a latência reduzida do modelo espectral e a especificidade elevada do modelo físico. Um \textit{ensemble} com votação ponderada é uma alternativa mais simples, porém com menor interpretabilidade.

\section{Conclusão}

O trabalho demonstrou a viabilidade de detectar falha de motor em VANTs de asa fixa com latência mediana inferior a 4\,s usando aprendizado não supervisionado sobre \textit{features} físicas e espectrais. O \texttt{all\_features} é o mais confiável operacionalmente (precisão 0{,}983, 69{,}6\,\% de detecção); o \texttt{fft\_features} é mais veloz e atua como precursor (latência mediana 2{,}64\,s, apenas 30\,\% de falsos positivos em voos normais). Nenhum modelo isolado satisfaz simultaneamente os requisitos de baixa latência e alta especificidade. Como trabalho futuro, recomenda-se implementar a arquitetura em cascata proposta, ajustar o \texttt{contamination} por validação temporal e ampliar o conjunto de dados com novos voos para cobrir cenários de falha não representados.

{\footnotesize
\bibliographystyle{plain}
\bibliography{refs}
}

\end{document}
